\section{连续运动可能需要固定量的新材料}\label{sec:density}
连续运动似乎是组织各个方向的一种自然方式。下一个定理说明，即使允许新单位针移到任意远处，这样组织方向仍可能迫使材料改变一个确定的量。

\begin{theorem}[短轴障碍]\label{thm:short-axis}
设 $K$ 是紧致原点星形体，在每个射影方向上都包含一条完整单位线段。假设 $K\subset B_{R_0}(0)$，其中 $R_0<1$。设 $m\ge3$，并假设对每个 $\ell=0,\ldots,m-1$，过原点的网格直线 $\mathbb R u_{\ell\pi/m}$ 满足 $K\cap\mathbb R u_{\ell\pi/m}\subset B_a(0)$，其中
\[
 a<\min(1-R_0,1/2).
\]
存在 $\varepsilon_K>0$，使每个包含连续闭单位针截面的紧致原点星形体 $Q$ 都满足
\[
 |Q\setminus K|\ge\varepsilon_K.
\]
$Q$ 的半径不受限制。
\end{theorem}

\begin{proof}
选取 $a<\rho_0<\rho_1<1/2$，并令 $\ell_0=1/2-\rho_1>0$。由紧致性和短轴假设，在每条射影直线 $q_\ell$ 周围存在一个角带 $V_\ell$，其中 $K$ 没有半径至少为 $\rho_0$ 的点。这个角带包含该直线的两条实际射线。$K$ 中平行于 $q_\ell$ 的完整单位线段不可能有零高度，因为过原点直线上的整个截面长度至多为 $2a<1$。

$K$ 中实际完整单位线段的集合是紧的。因此，存在以 $q_\ell$ 为内点的非退化闭区间 $J_\ell$ 及 $\eta_\ell>0$，使 $K$ 中方向落在 $J_\ell$ 内的每条实际单位线段都满足 $|h|>3\eta_\ell$。这个高度间隙只用于网格附近的方向窗；其他方向仍可有实际零高度单位线段。将 $J_\ell$ 缩入 $V_\ell$ 内，并进一步缩小 $\eta_\ell$，使得
\[
 q\in J_\ell,\quad |x|\ge\rho_1,\quad |h|\le\eta_\ell
 \quad\Longrightarrow\quad [xu_q+hn_q]\in V_\ell.
\]
这里方括号表示射影射线方向。由 $|h/x|\le\eta_\ell/\rho_1$，再利用 $J_\ell$ 与 $V_\ell$ 的补集之间的角度余量，就能一致地得到这个蕴含。

我们先对一整个低高度单位针区间收取面积。假设 $Q$ 的一个连续单位针截面在整个 $J_\ell$ 上满足 $|h(q)|<\eta_\ell$。将它的单位针写成 $(b(q)+[-1/2,1/2])u_q+h(q)n_q$。当 $b\ge0$ 时，保留正向外端长度为 $\ell_0$ 的尾段；当 $b\le0$ 时，保留负向外端同样长度的尾段。在 $b=0$ 时两段都保留。所选尾段的每一点，其纵向坐标的绝对值都至少为 $\rho_1$，其射线也都落在 $V_\ell$ 中。这些尾段形成一个紧致来源图；它们与原点张成的三角形的并集记为 $S$，包含于 $Q$。

$S$ 占据的每条射线都延伸到至少 $\rho_1$ 的半径。由于 $K$ 在 $V_\ell$ 中超出 $\rho_0$ 后没有材料，极坐标积分给出
\begin{equation}\label{eq:tail-relative-bank}
 |Q\setminus K|\ge |S\setminus B_{\rho_0}|
 \ge \left(1-\frac{\rho_0^2}{\rho_1^2}\right)|S|.
\end{equation}
非退化区间 $J_\ell$ 的有限多个旋转副本覆盖射影圆。若副本数为 $k_\ell$，则相应的 $S/\ell_0$ 的旋转副本一起组成一个原点星形体，在每个方向上都包含一条完整单位线段。对这个并集应用普适的 $1/10$ 定理，再用测度的次可加性，得到
\[
 |S|\ge\frac{\ell_0^2}{10k_\ell},\qquad
 |Q\setminus K|\ge
 \gamma_\ell:=\left(1-\frac{\rho_0^2}{\rho_1^2}\right)
 \frac{\ell_0^2}{10k_\ell}>0.
\]
这里估计的是一整个方向区间所需的面积。一个遥远而细薄的三角形可以有很小的面积；方向区间却迫使足够多的完整尾段出现，使下界定理能够应用。

接着把剩余的零代价配置放进一个紧空间。对单位线段 $N(q,b,h)$，令
\[
 f(q,b,h)=|\operatorname{conv}(0,N(q,b,h))\setminus K|.
\]
限制 $|h|\le H:=2(|K|+1)$，并在 $b\in\mathbb{R}$ 的两端加入无穷端点，将其紧致化。采用粘合 $(\pi,b,h)\sim(0,-b,-h)$，无穷端点也按此粘合。所得参数空间 $X$ 紧致且可度量。

在有限 $b$ 处，三角形按对称差面积连续变化，包括 $h=0$ 的情形，因此 $f$ 连续。当 $|b|$ 很大时，三角形内落在 $B_{R_0}$ 中的每一点，其收缩参数至多为 $R_0/(|b|-1/2)$。因此
\[
 |\operatorname{conv}(0,N)\cap K|
 \le \left(\frac{R_0}{|b|-1/2}\right)^2\frac{|h|}{2}.
\]
对有界的 $h$，右侧一致趋于零。在两个无穷端点处，将 $f$ 延拓为 $|h|/2$。

它的零点集由 $K$ 中实际包含的完整单位线段和所有零高度配置组成，包括无穷处的零高度配置。要看清这一点，注意新增面积为零的非退化三角形必须包含于闭集 $K$：三角形中任何落在 $K$ 外的点都有一个邻域，与该三角形的交集面积为正。相比之下，零高度三角形即使其单位线段不在 $K$ 内，面积仍为零。这些额外的径向配置说明，单针代价本身不足以证明定理。

在紧集 $q\in J_\ell$、$\eta_\ell\le|h|\le2\eta_\ell$ 上，函数 $f$ 有正的最小值 $\beta_\ell$。有限处的零代价状态必须是实际单位线段，但高度间隙排除了它们；无穷处状态的代价则为 $|h|/2>0$。现在假设
\[
 |Q\setminus K|<\min(1,\beta_0',\ldots,\beta_{m-1}',
                                \gamma_0,\ldots,\gamma_{m-1}),
 \qquad \beta_\ell'=\beta_\ell.
\]
该截面每个三角形的代价都至多为 $|Q\setminus K|$。此外，$|Q|\le |K|+1$，所以 $|h|/2\le |Q|$ 保证整个截面落在 $X$ 内。在每个连通的 $J_\ell$ 上，由连续性，$|h|$ 不能进入 $[\eta_\ell,2\eta_\ell]$。它必须始终小于 $\eta_\ell$，或者始终大于 $2\eta_\ell$。尾段支付排除了第一种情况。

对每个 $\ell$，从 $X$ 中删去开集
\[
 \{q\in\operatorname{int}J_\ell,\ |h|<2\eta_\ell\},
\]
将剩下的闭空间记为 $Y$。每个新增面积足够小的截面都落在 $Y$ 中。这个剩余空间中的零代价状态是 $K$ 中的实际单位线段，以及只存在于网格方向以外的径向状态。

网格射线将去掉原点的平面分成 $2m$ 个开扇区。$K$ 中实际单位线段的每个端点的模至少为 $1-R_0>a$，因此避开每条网格线。它的两个端点位于不同扇区：一个张角至多为 $\pi/3$、截在半径 $R_0<1$ 内的扇区，其直径小于一。用两个端点所在扇区的无序对给实际单位线段标号。这个标签局部常值。这里两个扇区只须不同；对踵条件将在径向状态中使用。

给剩余的径向零代价状态赋予一个对踵扇区对，分别包含 $u_q$ 和 $-u_q$；这个标签与 $b$ 无关，也包括 $b=\pm\infty$。由于网格方向已被删去，径向标签是连续的。在实际零高度单位线段处，两种标签相同：在 $B_{R_0}$ 内，过原点直线上的单位线段必须跨过原点。实际部分与径向部分都是闭的，因此它们的标签粘合为 $f|_Y$ 的整个零点集上的连续有限值标签。

对每个出现的扇区对 $\alpha$，选取一条严格分开这两个扇区的网格线。没有标签为 $\alpha$ 的实际单位线段平行于这条线，因为其端点的法向坐标异号。剩余的径向零代价状态也没有这个方向。因此，每个紧致标签纤维都缺少一个射影方向。可在 $Y$ 中为这有限多个纤维选取两两不交的开邻域 $U_\alpha$，且每个邻域仍缺少其指定方向。

这些邻域并集的补集是紧的；若非空，则 $f$ 在其上有严格为正的最小值，记为 $\beta_*>0$；若补集为空，任取 $\beta_*>0$。选取 $\varepsilon_K$，使其小于 $1$、所有 $\beta_\ell,\gamma_\ell$ 以及 $\beta_*$。一个新增面积小于 $\varepsilon_K$ 的截面，其连通像会包含在 $U_\alpha$ 的不交并中，因而包含在其中一个邻域内。这个邻域缺少一个方向，而截面在每个方向上都有一条单位针。矛盾证明了结论。
\end{proof}

短轴有限构造为定理~\ref{thm:short-axis} 提供了一个固定且具有竞争力的输入。对这个输入，更强的稠密性要求 $|Q\triangle K|\to0$ 也不成立。“固定”很重要：正的常数是在固定 $K$ 及其角向空隙之后才选出的。

等下确界问题仍须按它本来的形式来问。对任意候选替换，
\[
 |Q|-|K|=|Q\setminus K|-|K\setminus Q|.
\]
一个具有连续闭单位针截面的体，撤去的旧材料可能不少于新增材料。因此，这个障碍排除了几乎保留材料的连续实现，却没有比较所有可能重构的净面积。这一区别不是由一致半径上限造成的。证明将纵向逃逸紧致化，再用方向窗口上的集体面积下界排除相关的径向零代价状态。


这种区别在上述构造中确实发生。相干构造的上极限不超过 $c_{\mathrm{hl}}-\Delta_{\mathrm{coh}}$，而短轴族的面积趋于更大的 $c_{\mathrm{height}}$。因此可选出一个固定的原有限相干体 $G$，其面积小于所有充分晚的短轴输入。它不必包含那些输入；新增材料下界与更低的总面积可以同时成立。

